\chapter{Boolean Analysis --- Two Point}
%%% generated by the blueprint pipeline from TCSlib.BooleanAnalysis.Hypercontractivity.Cube.General.TwoPoint
\section{Overview}
Symmetric real-power inequalities establish the one-bit hypercontractive bound for finite exponents between one and two. The statements include the normalized scalar comparison and its affine-function consequence.

\section{Declarations}

\begin{lemma}[Symmetric power average is at least one]
\lean{BooleanAnalysis.Hypercontractivity.avg_rpow_ge_one}
\leanok
\difficulty{4}
\label{BooleanAnalysis.Hypercontractivity.avg_rpow_ge_one}
Let $p \ge 1$ be a real number and let $b \in [0,1]$. Then
\[
  1 \;\le\; \frac{(1+b)^p + (1-b)^p}{2}.
\]
\end{lemma}

\begin{lemma}[Monotonicity of symmetric convex sums]
\lean{BooleanAnalysis.Hypercontractivity.convex_sym_sum_mono}
\leanok
\difficulty{5}
\label{BooleanAnalysis.Hypercontractivity.convex_sym_sum_mono}
Let $f\colon\bbr\to\bbr$ be convex on the half-line $[0,\infty)$, and let $x,y$ be real
numbers with $0 \le x \le y \le 1$. Then
\[
  f(1+x) + f(1-x) \;\le\; f(1+y) + f(1-y).
\]
\end{lemma}

\begin{lemma}[Antitonicity in the exponent for nonpositive powers]
\lean{BooleanAnalysis.Hypercontractivity.rpow_sum_antitone_exponent}
\leanok
\difficulty{5}
\label{BooleanAnalysis.Hypercontractivity.rpow_sum_antitone_exponent}
Let $0 < x < 1$, and let $p \le q \le 0$ be real exponents. Then
\[
  (1+x)^p + (1-x)^p \;\ge\; (1+x)^q + (1-x)^q.
\]
\end{lemma}

\begin{lemma}[A two-point convexity inequality]
\lean{BooleanAnalysis.Hypercontractivity.h_alpha_ineq}
\leanok
\uses{BooleanAnalysis.Hypercontractivity.rpow_sum_antitone_exponent}
\difficulty{7}
\label{BooleanAnalysis.Hypercontractivity.h_alpha_ineq}
Let $r,s\in\bbr$ satisfy $0\le r\le s\le1$, let $t\in[0,1]$, and set
$c=\sqrt{r/s}$ when $s>0$ and $c=0$ when $s=0$. With $0^0=1$, one has
\[
  (1+t)^r-(1-t)^r\ge c\bigl((1+ct)^s-(1-ct)^s\bigr).
\]
\end{lemma}

\begin{lemma}[Integrated two-point hypercontractive inequality]
\lean{BooleanAnalysis.Hypercontractivity.integrated_h_alpha_ineq}
\leanok
\uses{BooleanAnalysis.Hypercontractivity.h_alpha_ineq}
\difficulty{6}
\label{BooleanAnalysis.Hypercontractivity.integrated_h_alpha_ineq}
Let $p,q,b\in\bbr$ satisfy $1\le p\le q\le2$ and $0\le b\le1$.
Set $\rho=\sqrt{(p-1)/(q-1)}$ when $q>1$ and $\rho=0$ when $q=1$. Then
\[
  \frac{(1+\rho b)^q+(1-\rho b)^q-2}{q}
  \le\frac{(1+b)^p+(1-b)^p-2}{p}.
\]
\end{lemma}

\begin{lemma}[Tangent line inequality for $x^r$ at $x=1$]
\lean{BooleanAnalysis.Hypercontractivity.rpow_ge_one_add_mul_sub}
\leanok
\difficulty{4}
\label{BooleanAnalysis.Hypercontractivity.rpow_ge_one_add_mul_sub}
For every real number $x \ge 0$ and every exponent $r \ge 1$, the power $x^r$ lies above
its tangent line at $x = 1$; that is,
\[
x^r \ge 1 + r\,(x - 1).
\]
\end{lemma}

\begin{theorem}[General two-point inequality, unit case]
\lean{BooleanAnalysis.Hypercontractivity.two_point_ineq_general_unit}
\leanok
\uses{BooleanAnalysis.Hypercontractivity.integrated_h_alpha_ineq, BooleanAnalysis.Hypercontractivity.rpow_ge_one_add_mul_sub}
\difficulty{4}
\label{BooleanAnalysis.Hypercontractivity.two_point_ineq_general_unit}
\statementsource{boolean-ch09-hypercontractivity-kkl}{pdf-4ba63c9394e5-p001-b002}
\statementsource{boolean-ch10-advanced-hypercontractivity}{
  pdf-9d1751394783-p005-b001,
  pdf-9d1751394783-p005-b002
}
Let $p,q,b\in\bbr$ satisfy $1\le p\le q\le2$ and $0\le b\le1$.
Set $\rho=\sqrt{(p-1)/(q-1)}$ when $q>1$ and $\rho=0$ when $q=1$. Then
\[
  \frac{(1+\rho b)^q+(1-\rho b)^q}{2}
  \le\left(\frac{(1+b)^p+(1-b)^p}{2}\right)^{q/p}.
\]
\end{theorem}

\begin{theorem}[One-bit low-norms hypercontractivity]
\lean{BooleanAnalysis.Hypercontractivity.low_norms_one_bit}
\leanok
\uses{BooleanAnalysis.BoolCube, BooleanAnalysis.BooleanFunc, BooleanAnalysis.Hypercontractivity.expect_abs_rpow_one_bit, BooleanAnalysis.Hypercontractivity.two_point_ineq_general_unit, BooleanAnalysis.boolToSign, BooleanAnalysis.chiS, BooleanAnalysis.expect, BooleanAnalysis.fourierCoeff, BooleanAnalysis.innerProduct, BooleanAnalysis.noiseOp, BooleanAnalysis.uniformWeight}
\difficulty{7}
\label{BooleanAnalysis.Hypercontractivity.low_norms_one_bit}
\statementsource{boolean-ch10-advanced-hypercontractivity}{
  pdf-9d1751394783-p004-b006,
  pdf-9d1751394783-p004-b007,
  pdf-9d1751394783-p005-b001,
  pdf-9d1751394783-p005-b002
}
\statementsource{boolean-ch09-hypercontractivity-kkl}{pdf-4ba63c9394e5-p001-b002}
Let $1 < p \le q \le 2$, let $f : \{0,1\}^1 \to \bbr$ be a Boolean function of arity
$1$, and set $\rho = \sqrt{(p-1)/(q-1)}$. Then, with expectations taken under the
uniform measure and $T_\rho$ the noise operator,
\[
\bigl(\E\bigl[\,\abs{T_\rho f}^q\,\bigr]\bigr)^{1/q} \;\le\;
\bigl(\E\bigl[\,\abs{f}^p\,\bigr]\bigr)^{1/p}.
\] For $d\in\mathbb N$, $h:\{0,1\}^d\to\bbr$, and $\tau\in\bbr$, the spectral noise operator is $T_\tau h(x)=\sum_{S\subseteq[d]}\tau^{|S|}\widehat h(S)\chi_S(x)$, where $\chi_S(x)=(-1)^{\sum_{i\in S}x_i}$ and $\widehat h(S)=\E[h\chi_S]$ under the uniform measure.
\end{theorem}
